\chapter{Frontend Implementation Detail}
\label{chap:features}

The frontend is the core of the current codebase. It is a component-based application managed by React v17.0.2.

\section{Visual Component Strategy (Material UI \& Leaflet)}

The User Interface relies on \texttt{@material-ui/core} (v4) for layout and \texttt{react-leaflet} for geospatial rendering.

\subsection{Map Rendering Engine (Maps.js)}

The map is not a static image but an interactive Canvas/SVG layer managed by Leaflet.

\textbf{Tile Layer Architecture:}
\begin{itemize}
    \item \textbf{Tile Source}: Fetches graphical map tiles from OpenStreetMap
    \begin{itemize}
        \item Default: \texttt{https://\{s\}.tile.openstreetmap.org/\{z\}/\{x\}/\{y\}.png}
    \end{itemize}
\end{itemize}

\textbf{Vector Layers:}
\begin{itemize}
    \item To draw the routes, the app expects an array of \texttt{[lat, long]} pairs from the backend
    \item These are passed to the \texttt{<Polyline />} component, which draws the colored path overlays
    \begin{itemize}
        \item Blue: Car routes
        \item Green: Bike routes
        \item Orange: Walking routes
    \end{itemize}
\end{itemize}

\textbf{Camera Markers:}
\begin{itemize}
    \item Custom camera icons rendered using Leaflet markers
    \item Click handler triggers popup with real-time camera image
    \item Auto-zoom functionality when camera is selected
\end{itemize}

\section{Routing \& Navigation (App.js)}

The application uses React Router v6 to create a seamless SPA experience. The browser does not reload when switching views; instead, the Virtual DOM is updated.

\begin{table}[htbp]
\centering
\begin{tabular}{|l|l|p{6cm}|}
\hline
\textbf{Route Path} & \textbf{Component} & \textbf{Description} \\
\hline
\texttt{/} or \texttt{/routes} & RoutesPage & Primary dashboard for route planning \\
\hline
\texttt{/busmap} & BusMapPage & Specialized view for bus network \\
\hline
\texttt{/cameras} & CameraMap & Traffic camera viewing and search \\
\hline
\texttt{/history} & HistoryPage & User's saved locations and routes \\
\hline
\end{tabular}
\caption{Application routing structure}
\label{tab:routes}
\end{table}

\section{State Management}

State is currently handled locally within components using React Hooks (\texttt{useState}, \texttt{useEffect}).

\textbf{Input State:}
\begin{itemize}
    \item SearchBox captures user keystrokes
    \item Debounced search (500ms) to optimize API calls
\end{itemize}

\textbf{Data State:}
\begin{itemize}
    \item \texttt{SearchBoxRoutes} holds the JSON object returned by the server
    \item \texttt{filteredCameras} maintains filtered camera list
    \item \texttt{selectedLocation} tracks current route start point
    \item \texttt{selectPosition} tracks current route end point
\end{itemize}

\textbf{Persistence State:}
\begin{itemize}
    \item \texttt{recentHistory} stores user's saved locations and routes
    \item Loaded from backend on component mount
\end{itemize}

\section{Component Hierarchy}

The application follows a hierarchical component structure:

{\footnotesize
\begin{forest}
  for tree={
    font=\normalfont,
    grow'=0,
    child anchor=west,
    parent anchor=south,
    anchor=west,
    calign=first,
    edge path={
      \noexpand\path [draw, \forestoption{edge}]
      (!u.south west) +(5pt,0) |- node[fill,inner sep=0.8pt] {} (.child anchor)\forestoption{edge label};
    },
    before typesetting nodes={
      if n=1
        {insert before={[,phantom]}}
        {}
    },
    fit=band,
    s sep=1.5pt,
    l sep=6pt,
  }
  [App.js
    [Header\_v2.js (Nav)]
    [Routes
      [RoutesPage.js
        [Maps.js
          [MapLayerControl.js
            [OSMVector.js]
          ]
          [Marker (Start)]
          [Marker (End)]
          [Polyline]
        ]
        [SearchBoxRoutes.js]
        [History Panel]
      ]
      [CameraMap.js
        [Maps.js]
        [Search Bar]
        [Camera Markers]
        [CameraPopup.js]
      ]
      [BusMapPage.js]
      [HistoryPage.js]
    ]
  ]
\end{forest}
}